
{\actuality}
 Развитие робототехники ставит перед исследователями различные задачи навигации автономных роботизированных средств.
Принцип автономной навигации мобильных роботов включает в себя поиск наилучшего маршрута, т.е. оптимального пути \cite{LiuWang2023}.
Самый простой способ - это перемещение между известными ориентирами по кратчайшему расстоянию.
Для оптимизации этого способа, необходимо учитывать начальное и конечное состояния.
Под оптимизацией понимается тот факт, что алгоритмы минимизируют одну или несколько целевых функций.
Однако, например, в случае роботов, которые участвуют в операциях по поиску и спасению людей, целевая функция определяется приоритетами, установленными оператором \cite{Calisi2007}.
По-другому определяется функция картографирования больших пространств, она может быть определена посредством снятия целевых данных датчиками с окружения.
Дополнительные ограничения накладывают характеристики среды, влияющие на продвижение робота.
Так, расчет движения робота в стоячей воде, в отличии от движения по плоской твердой поверхности, требует учитывать инерцию среды, а в расчетах появляются параметры отражающие дополнительные степени свободы.

В области управления судами есть несколько основополагающих работ \cite{Alexandersson2023}, которые задали тренд в моделировании движения судов в открытой воде с учетом течений и без них.
Первой такой работой стала модель Абковица, описанная в 1964 году \cite{Abkowitz1964}, в которой он представил комплексную модель маневрирования судна, также известную, как модель типа Абковица, рассматривая корпус, гребной винт и руль направления как целостную систему.
Позже Огава и Касаи (1978)\cite{ogawa1978} представили математическую модель маневрирования, более известную в международной литературе как Manoeuvring Mathematical Group Model(MMG model) для представления гидродинамики корпуса, гребного винта, рулевой системы и их взаимодействия по отдельности.
В этой модели гидродинамика судна характеризуется гидродинамическими производными, которые обычно получаются путем проведения испытаний на реальной модели.
В последствии MMG модель и модель Абковица неоднократно тестировались на реальных судах.

Серьезный вклад в области систем управления движением морских судов, гидродинамики, теории управления, систем наведения и навигации в данный момент сделал профессор Фоссен\cite{fossen2021handbook}.
Современные гидродинамические модели, встречающиеся в литературе, описаны в матрично-векторном виде предложенном Фоссеном (1991, 1994, 2011)\cite{fossen1991}\cite{fossen1994}\cite{fossen2011handbook}.
Матрично-векторный вид записи удобен при проектировании систем управления, поскольку в расчетах можно использовать такие свойства, как симметрия, кососимметрия и положительная определенность.
Математические модели, описанные Фоссеном, являются стандартом проектирования систем управления движением морских судов.
Они подробно описаны популярным языком программирования и являются открытым программным обеспечением.
Все это позволяет применить модель Фоссена для множества смежных исследовательских задач.

Данная работа является продолжением работ профессора Матвеева о задачах навигации неголономных роботов (2016, 2023)\cite{matveev2016}\cite{matveev2023}.
В своих статьях, профессор Матвеев исследует задачи навигации неголономных роботов типа Дубинса.
Его закон управления роботом для нахождения границ скалярного поля, можно считать универсальным для разного типа роботов (морских, воздушных, сухопутных).
Закон управления был опробован на физических моделях сухопутного робота типа Дуббинса\cite{semakova2017phdthesis} и на морском роботе катамаранного типа\cite{matveev2023}.
Это значит, что можно применить закон профессора Матвеева для управления судном катамаранного типа описанным в численной модели профессора Фоссена в задаче поиска изолинии в скалярном поле.
В работах Матвеева точная природа скалярного поля не уточняется, оно рассматривается в общем виде.
Однако, актуальность решения задачи поиска изолинии в скалярном поле может быть сопоставлена с реальной задачей экологического мониторинга при разливе нефти в океане.
Определение границ нефтяного пятна позволяет вычислить его размер, скорость распространения, объем нефти по значению концентрации в данной точке изолинии и другие параметры необходимые для ликвидации экологической угрозы \cite{Aznar2014}.

Цель исследования, перенести закон управления неголономным роботом на качественную апробированную открытую программную модель морского судна и проверить его работу.
В основные задачи исследовательской работы входит:
\begin{enumerate}
    \item Поиск открытой и апробированной компьютерной модели неголономного судна, схожего по принципу движения с роботом типа Дубинса.
    \item  Поиск способов описания и задания скалярного поля, для навигации в нем.
    \item Модернизация компьютерной модели профессора Фоссена управляемого судна катамаранного типа.
    \item Оценка точности работы закона управления в компьютерной модели.
\end{enumerate}

\ifsynopsis
Этот абзац появляется только в~автореферате.
Для формирования блоков, которые будут обрабатываться только в~автореферате,
заведена проверка условия \verb!\!\verb!ifsynopsis!.
Значение условия задаётся в~основном файле документа (\verb!synopsis.tex! для
автореферата).
\else
Этот абзац появляется только в~диссертации.
Через проверку условия \verb!\!\verb!ifsynopsis!, задаваемого в~основном файле
документа (\verb!dissertation.tex! для диссертации), можно сделать новую
команду, обеспечивающую появление цитаты в~диссертации, но~не~в~автореферате.
\fi

% {\progress}
% Этот раздел должен быть отдельным структурным элементом по
% ГОСТ, но он, как правило, включается в описание актуальности
% темы. Нужен он отдельным структурынм элемементом или нет ---
% смотрите другие диссертации вашего совета, скорее всего не нужен.

{\aim} данной работы является \ldots

Для~достижения поставленной цели необходимо было решить следующие {\tasks}:
\begin{enumerate}[beginpenalty=10000] % https://tex.stackexchange.com/a/476052/104425
  \item Исследовать, разработать, вычислить и~т.\:д. и~т.\:п.
  \item Исследовать, разработать, вычислить и~т.\:д. и~т.\:п.
  \item Исследовать, разработать, вычислить и~т.\:д. и~т.\:п.
  \item Исследовать, разработать, вычислить и~т.\:д. и~т.\:п.
\end{enumerate}


{\novelty}
\begin{enumerate}[beginpenalty=10000] % https://tex.stackexchange.com/a/476052/104425
  \item Впервые \ldots
  \item Впервые \ldots
  \item Было выполнено оригинальное исследование \ldots
\end{enumerate}

{\influence} \ldots

{\methods} \ldots

{\defpositions}
\begin{enumerate}[beginpenalty=10000] % https://tex.stackexchange.com/a/476052/104425
  \item Первое положение
  \item Второе положение
  \item Третье положение
  \item Четвертое положение
\end{enumerate}

{\reliability} полученных результатов обеспечивается \ldots \ Результаты находятся в соответствии с результатами, полученными другими авторами.


{\probation}
Основные результаты работы докладывались~на:
перечисление основных конференций, симпозиумов и~т.\:п.

{\contribution} Автор принимал активное участие \ldots

\ifnumequal{\value{bibliosel}}{0}
{%%% Встроенная реализация с загрузкой файла через движок bibtex8. (При желании, внутри можно использовать обычные ссылки, наподобие `\cite{vakbib1,vakbib2}`).
    {\publications} Основные результаты по теме диссертации изложены
    в~XX~печатных изданиях,
    X из которых изданы в журналах, рекомендованных ВАК,
    X "--- в тезисах докладов.
}%
{%%% Реализация пакетом biblatex через движок biber
    \begin{refsection}[bl-author, bl-registered]
        % Это refsection=1.
        % Процитированные здесь работы:
        %  * подсчитываются, для автоматического составления фразы "Основные результаты ..."
        %  * попадают в авторскую библиографию, при usefootcite==0 и стиле `\insertbiblioauthor` или `\insertbiblioauthorgrouped`
        %  * нумеруются там в зависимости от порядка команд `\printbibliography` в этом разделе.
        %  * при использовании `\insertbiblioauthorgrouped`, порядок команд `\printbibliography` в нём должен быть тем же (см. biblio/biblatex.tex)
        %
        % Невидимый библиографический список для подсчёта количества публикаций:
        \phantom{\printbibliography[heading=nobibheading, section=1, env=countauthorvak,          keyword=biblioauthorvak]%
        \printbibliography[heading=nobibheading, section=1, env=countauthorwos,          keyword=biblioauthorwos]%
        \printbibliography[heading=nobibheading, section=1, env=countauthorscopus,       keyword=biblioauthorscopus]%
        \printbibliography[heading=nobibheading, section=1, env=countauthorconf,         keyword=biblioauthorconf]%
        \printbibliography[heading=nobibheading, section=1, env=countauthorother,        keyword=biblioauthorother]%
        \printbibliography[heading=nobibheading, section=1, env=countregistered,         keyword=biblioregistered]%
        \printbibliography[heading=nobibheading, section=1, env=countauthorpatent,       keyword=biblioauthorpatent]%
        \printbibliography[heading=nobibheading, section=1, env=countauthorprogram,      keyword=biblioauthorprogram]%
        \printbibliography[heading=nobibheading, section=1, env=countauthor,             keyword=biblioauthor]%
        \printbibliography[heading=nobibheading, section=1, env=countauthorvakscopuswos, filter=vakscopuswos]%
        \printbibliography[heading=nobibheading, section=1, env=countauthorscopuswos,    filter=scopuswos]}%
        %
        \nocite{*}%
        %
        {\publications} Основные результаты по теме диссертации изложены в~\arabic{citeauthor}~печатных изданиях,
        \arabic{citeauthorvak} из которых изданы в журналах, рекомендованных ВАК%
        \ifnum \value{citeauthorscopuswos}>0%
            , \arabic{citeauthorscopuswos} "--- в~периодических научных журналах, индексируемых Web of~Science и Scopus%
        \fi%
        \ifnum \value{citeauthorconf}>0%
            , \arabic{citeauthorconf} "--- в~тезисах докладов.
        \else%
            .
        \fi%
        \ifnum \value{citeregistered}=1%
            \ifnum \value{citeauthorpatent}=1%
                Зарегистрирован \arabic{citeauthorpatent} патент.
            \fi%
            \ifnum \value{citeauthorprogram}=1%
                Зарегистрирована \arabic{citeauthorprogram} программа для ЭВМ.
            \fi%
        \fi%
        \ifnum \value{citeregistered}>1%
            Зарегистрированы\ %
            \ifnum \value{citeauthorpatent}>0%
            \formbytotal{citeauthorpatent}{патент}{}{а}{}%
            \ifnum \value{citeauthorprogram}=0 . \else \ и~\fi%
            \fi%
            \ifnum \value{citeauthorprogram}>0%
            \formbytotal{citeauthorprogram}{программ}{а}{ы}{} для ЭВМ.
            \fi%
        \fi%
        % К публикациям, в которых излагаются основные научные результаты диссертации на соискание учёной
        % степени, в рецензируемых изданиях приравниваются патенты на изобретения, патенты (свидетельства) на
        % полезную модель, патенты на промышленный образец, патенты на селекционные достижения, свидетельства
        % на программу для электронных вычислительных машин, базу данных, топологию интегральных микросхем,
        % зарегистрированные в установленном порядке.(в ред. Постановления Правительства РФ от 21.04.2016 N 335)
    \end{refsection}%
    \begin{refsection}[bl-author, bl-registered]
        % Это refsection=2.
        % Процитированные здесь работы:
        %  * попадают в авторскую библиографию, при usefootcite==0 и стиле `\insertbiblioauthorimportant`.
        %  * ни на что не влияют в противном случае
        \nocite{vakbib2}%vak
        \nocite{patbib1}%patent
        \nocite{progbib1}%program
        \nocite{bib1}%other
        \nocite{confbib1}%conf
    \end{refsection}%
        %
        % Всё, что вне этих двух refsection, это refsection=0,
        %  * для диссертации - это нормальные ссылки, попадающие в обычную библиографию
        %  * для автореферата:
        %     * при usefootcite==0, ссылка корректно сработает только для источника из `external.bib`. Для своих работ --- напечатает "[0]" (и даже Warning не вылезет).
        %     * при usefootcite==1, ссылка сработает нормально. В авторской библиографии будут только процитированные в refsection=0 работы.
}

